В данной главе проводится детальный анализ результатов экспериментального сравнения различных подходов к извлечению библиографии, обсуждаются качественные характеристики ошибок, оценивается вычислительная эффективность методов и описывается процесс интеграции разработанного решения в производственную среду информационной системы ИСАНД.

\section{Сравнительный анализ подходов}
На этапах локализации границ библиографических блоков и построчной сегментации записей инструмент GROBID показал наилучшие результаты при обработке англоязычного корпуса. Высокое качество парсинга обусловлено спецификой архитектуры данного инструмента: в процессе предварительного разбора PDF-файла алгоритм извлекает не только символьные цепочки, но и пространственно-геометрические параметры макета страницы. Наличие этих данных позволяет модели стабильно дифференцировать текстовые колонки, элементы колонтитулов и тело основной статьи.

Напротив, схема на базе регулярных выражений и градиентного бустинга оказалась уязвимой перед стилевым многообразием вёрстки. Детерминированные правила успешно распознают явную нумерацию вроде «<<[1]>>» или «<<[2]>>», но теряют эффективность при анализе монолитных абзацев, где ссылки разделяются обычными точками с пробелами. Сбои фиксируются и при обработке строк с выступами, если первичный программный экстрактор исказил координатную сетку PDF-слоя.

Экспериментальная проверка БЯМ на исходных текстах статей показала их слабую применимость к подзадаче $g_1$. Ограниченность контекста и размытие весов внимания на длинных цепочках токенов приводят к ошибкам детекции начала списка литературы. Наконец, из-за генеративных свойств БЯМ возникает эффект скрытого редактирования. Происходит отбрасывание строк или автокоррекция опечаток, что разрушает логику точной сегментации, где критически важен исходный синтаксис. Из этого следует, что для задач $g_1$ и $g_2$ использование детерминированных инструментов вроде GROBID остается базовым стандартом.

Картина принципиально меняется на этапе извлечения библиографических областей, где GROBID показал резкое падение качества при извлечении сложных сущностей, в то время как дообученная модель Llama-3-8B достигает высокого макро-$F_1 = 0,84$.

Низкое качество GROBID на этапе $g_3$ объясняется ограничениями CRF/BiLSTM архитектуры. GROBID решает задачу структурирования как последовательную разметку токенов в формате BIO: Begin, Inside, Outside. Условные случайные поля оперируют исключительно соседними токенами и жёсткими эвристическими признаками. Когда GROBID сталкивается с нестандартным оформлением авторов (например, инвертированный порядок «Красоткин Семён Александрович», отсутствие пробелов между инициалами «С.А.Красоткин», наличие иностранных частиц «фон», «де», «ла»), модель не может сопоставить текущий токен с глобальным смыслом записи. Точка, которая в одном стиле разделяет фамилию и инициалы, в другом стиле может завершать название журнала, что приводит к каскадной ошибке разметки и ошибочному отнесению фамилии к полю заголовка.

В свою очередь, БЯМ благодаря механизму self-attention обладают способностью улавливать глобальный семантический контекст всей записи целиком. Модель понимает, что токены, следующие сразу после начала строки и имеющие форму заглавных букв с точками, с высокой вероятностью относятся к авторам, даже если пунктуация нарушена или отсутствует. В текущей версии системы LoRA не используется, потому что эксперимент не удался из-за зашумлённых данных, а базовая модель и так даёт приемлемое качество $F_1=0,84$.

Для глубокого понимания пределов применимости разработанной системы проведён качественный анализ ошибок, классифицированный по источникам возникновения.

Наибольшее количество неустранимых ошибок связано с кумулятивным эффектом шумов на этапе извлечения текста (действие отображения $g_1$). Если библиотека PDFplumber или PyMuPDF объединяет конец одной строки с началом другой без пробела, GROBID формирует единый токен, который семантически не существует. 

Специфика библиографических данных, предоставленный из Math-Net.Ru, заключается в их синтаксической неоднородности, вызванной обилием математического аппарата в названиях публикаций. Хотя модель DeepSeek обнаружила высокую инвариантность к пунктуационным искажениям верстки, в ходе вычислительных экспериментов проявился систематический дефект. Алгоритм склонен опускать или деформировать математические выражения внутри заголовков, принудительно подменяя отсутствующие в словаре токенизатора спецсимволы на текстовые маркеры-заглушки. Данный сбой обусловлен фундаментальными особенностями авторегрессионных архитектур. Для преодоления этого барьера необходима интеграция дополнительного конвейера постпроцессинга, восстанавливающего исходные формулы с помощью регулярных масок.

Критическое значение для стабильности конвейера имеет дефект слияния независимых строк на этапе $g_2$. В единичных сценариях инструмент GROBID объединял две смежные библиографические записи, если в исходном PDF-документе отсутствовал межстрочный вертикальный интервал, а авторы отказались от сквозной нумерации списка. Передача подобного композитного токена на вход языковой модели приводила к системному сбою: Llama-3-8B пыталась агрегировать атрибуты двух разных публикаций в рамках единой JSON-схемы, что вызывало принудительное отсечение информации либо генерацию синтаксически нарушенного кода. Данный факт наглядно доказывает каскадный характер ошибок: погрешность на шаге сегментации $g_2$ полностью блокирует возможность корректного извлечения семантических полей.

Наряду с показателями полноты разбора, интеграция разработанных модулей в промышленный контур сдерживается требованиями к объёму видеопамяти и регламентированным временем отклика системы. Процедура отладки конвейера включала детальное сопоставление временных и ресурсных характеристик исследуемых подходов.

Инструмент GROBID, реализованный на языке Java и использующий компактные модели условных случайных полей, характеризуется низкими требованиями к вычислительным ресурсам и минимальным объёмом оперативной памяти. Это позволяет эффективно развертывать его в многопоточном режиме на стандартной серверной инфраструктуре общего назначения без привлечения специализированных графических ускорителей.

В свою очередь, развёртывание модели Llama-3-8B в исходной точности FP16 блокируется жёсткими лимитами видеопамяти. Проблема дефицита VRAM решена переходом к 4-битному сжатию весовых коэффициентов в формате GGUF $Q4\_K\_M$, сделав систему совместимой с доступными графическими ускорителями массового сегмента. Однако вычислительная сложность авторегрессионного декодирования накладывает ограничения на скорость обработки каждой отдельной ссылки. При анализе объёмных библиографических списков кумулятивная задержка языковой модели становится критической, что на порядок замедляет работу по сравнению с rule-based аналогами.

Архитектурный компромисс достигнут за счёт гибридизации системы, устранившей дисбаланс вычислительных скоростей. Первичная детекция и грубая сегментация возлагаются на быстродействующий детерминированный парсер, тогда как интеллектуальный семантический анализ изолированных текстовых строк выполняется большой языковой моделью в фоновом, асинхронном режиме.

\section{Внедрение в систему ИСАНД}

Разработанный конвейер интегрирован в качестве прототипа в информационную систему анализа научной деятельности (ИСАНД) в виде микросервисной архитектуры. В связи с разительным отличием в требованиях к вычислительным ресурсам, система была разделена на два независимых контура: контур первичной обработки и контур глубокого структурирования.

Интеграция разработанного модуля в информационную систему анализа научной деятельности (ИСАНД) осуществлялась по микросервисной архитектуре с использованием контейнеризации Docker.

На текущий момент в продуктивной среде ИСАНД развернут и постоянно функционирует модуль на основе GROBID. Сервис запущен в виде фонового Java-процесса, обеспечивающего REST API для приема PDF-файлов. По умолчанию все запросы направляются в grobid-worker.

Модуль на базе $Llama-3-8B-Instruct.Q4\_K\_M.gguf$ спроектирован как отдельный асинхронный Python-сервис с использованием FastAPI и библиотеки `llama-cpp-python`. На текущий момент он не развернут на общем веб-сервере ИСАНД из-за необходимости выделенного GPU-узла, но прошёл экспериментальное нагрузочное тестирование на нескольких выпусках журналов теории управления. 

Процесс взаимодействия двух модулей выглядит следующим образом:
\begin{enumerate}
    \item Пользователь загружает PDF в ИСАНД.
    \item Система отправляет файл в GROBID, получая обратно TEI/XML с сырыми записями.
    \item Сырые записи сохраняются в базу данных (обеспечивая мгновенный ответ пользователю).
    \item В фоновом режиме (через брокер сообщений) сырые строки передаются в Llama-сервис.
    \item Llama-сервис возвращает enrichment-данные (JSON с уточненными полями авторов, заголовков и т.д.), которые обновляются в базе данных ИСАНД.
\end{enumerate}
Такая двухуровневая архитектура позволяет системе не заставлять пользователя ждать 15 секунд при загрузке файла, немедленно показывая результат работы GROBID, и параллельно догружать качественную семантику с помощью БЯМ.

\section{Оценка достоверности полученных результатов}

Достоверность результатов в данной работе обеспечивается соблюдением стандартных практик машинного обучения, корректным разделением данных, обоснованным выбором метрик и фиксацией экспериментального окружения. Поскольку задача разбора библиографии чувствительна к шуму, стилевому разнообразию и артефактам препроцессинга, валидация строится на воспроизводимости условий и прозрачной оценке качества.

Обучающая, валидационная и тестовая выборки формировались на уровне целых документов, а не отдельных библиографических записей. Это исключает ситуацию, когда фрагменты одной статьи попадают в разные подвыборки, что могло бы искусственно завысить метрики за счёт контекстуального сходства. Все этапы предобработки (очистка текста, нормализация кодировок, фильтрация артефактов) применялись к каждой подвыборке независимо, без использования статистик или словарей, вычисленных по всему корпусу. Таким образом, риск утечки данных сведён к минимуму, а тестовая выборка имитирует реальные условия обработки ранее не встречавшихся публикаций.

Для оценки использовалась символьная $F_1$-мера в агрегации макро-среднего. Отказ от токеновой метрики обусловлен особенностями работы субсловных токенизаторов BPE/SentencePiece, которые по-разному сегментируют кириллические фамилии, инициалы и знаки препинания по сравнению с эталонной разметкой. Символьное сравнение напрямую отражает пригодность извлечённых полей для последующей нормализации и сопоставления с внешними базами. Макро-усреднение по полям выбрано для того, чтобы редкие, но важные элементы: соавторы, специфичные обозначения источников, DOI, -- не маскировались частотными полями вроде года издания или названия журнала.

Дообучение проводилось с использованием конфигурации LoRA-адаптеров, рекомендованной фреймворком Unsloth для моделей семейства Llama-3 ($r=16$, $\alpha=32$, \texttt{lora\_dropout}=0,05). Расширенный поиск гиперпараметров не проводился, поскольку основная цель эксперимента заключалась в проверке применимости подхода к русскоязычной библиографии, а не в поиске теоретического потолка качества. Для обеспечения воспроизводимости зафиксированы версии ключевых библиотек (\texttt{transformers}, \texttt{peft}, \texttt{unsloth}, \texttt{bitsandbytes}), параметры генерации (\texttt{temperature}=0,1, \texttt{top\_p}=0,9, \texttt{max\_new\_tokens}=512) и спецификация вычислительного окружения (CUDA 12.1, Python 3.10). Исходный код пайплайна и скрипты оценки версионированы.

Корпус Math-Net.Ru предоставлен на условиях, согласованных с правообладателем, и используется исключительно в исследовательских целях. Результаты валидны для текстовых PDF и LaTeX-исходников научных статей, где библиография оформлена отдельным разделом в конце документа. Система не оптимизирована для полностью растровых сканов низкого разрешения, постраничных сносок или документов с отсутствующим текстовым слоем, так как в этих случаях требуется отдельный этап OCR и восстановления логической структуры, выходящий за рамки текущего конвейера.

Ограниченность вычислительных ресурсов предопределила проведение одиночной итерации эксперимента без процедуры многократного усреднения по случайным инициализациям весов. Подобный подход обоснован высокой стоимостью выполнения современных больших языковых моделей, хотя полученные численные значения и остаются жестко привязанными к текущему разбиению данных. Вместе с тем, репрезентативность результатов обеспечивается за счет изоляции документов на уровне выборок, автономного препроцессинга, применения инвариантных к токенизации метрик и фиксации программно-аппаратной среды. Достигнутый уровень макро-\(F1 = 0{,}84\) на кириллическом материале доказывает состоятельность гибридной архитектуры и служит основанием для внедрения разработанного программного модуля в контур ИСАНД.

\section{Дальнейшие работы}

Факт достижения целевых метрик не исключает наличия у разработанного пайплайна архитектурных лимитов. Дальнейшая модернизация комплекса предполагает решение трех прикладных задач: оптимизация GROBID под специфику кириллических изданий, разбор распределенных по тексту библиографических сносок и внедрение фильтров для блокировки недостоверных генеративных данных.

Инструмент GROBID исторически ориентирован на англоязычный сегмент научной периодики, что приводит к падению качества на этапах $g_2$ и $g_3$ при обработке русскоязычных текстов. Это обусловлено следующими факторами: высокая вариативность антропонимов и специфические разделители.

Для решения этой задачи планируется полный цикл переобучения CRF-каскада GROBID на базе корпуса Math-Net.Ru. С помощью фреймворка \textit{grobid-trainer} PDF-файлы будут конвертированы в размеченный TEI-формат. Вектор признаков токена $x_{ij}$ будет расширен за счет кириллических словарей и регулярных выражений для детекции маркеров ГОСТ, что позволит повысить полноту извлечения до уровня англоязычных аналогов.

Текущая реализация опирается на допущение, что ссылки локализованы в едином блоке «Список литературы». Однако в гуманитарных дисциплинах доминируют постраничные сноски, также известные как footnotes. Их интеграция в конвейер требует перехода к анализу макета документа из-за следующих сложностей: сноски расположены внизу страницы и часто не отделены от основного текста жесткими маркерами; сноска может переноситься на следующую страницу. Кроме того, гуманитарные тексты активно используют сокращения: «там же», «Указ. соч.», которые бессмысленно парсить в изоляции без модуля разрешения кореферентности.

В качестве решения планируется внедрение визуально-текстовых моделей, например, класса LayoutLM или YOLO-Ref, для сегментации пространства страницы на зоны: (\textit{MainText}, \textit{Footer}, \textit{FootnoteRegion}). После извлечения зоны сноски токены будут направляться в парсер, оснащенный стековой памятью для связывания сокращенных ссылок с полными описаниями.

Использование БЯМ для написания статей привело к появлению синтаксически идеальных, но физически несуществующих библиографических записей. Поскольку структура таких «фантомных» ссылок безупречна, этапы $g_2$ и $g_3$ отработают без ошибок, и мусор попадет в базу данных. 

Для фильтрации требуется расширить математическую модель четвертым отображением верификации $g_4$, формируя композицию:
 $$ f_{\text{verified}} = g_4 \circ g_3 \circ g_2 \circ g_1
 $$ где $g_4 : C \to C \times \{0, 1\}$ — функция, сопоставляющая записи бинарный флаг истинности на основе внешних баз данных.

\begin{algorithm}[H]
\caption{Алгоритм верификации библиографической записи}
\label{alg:verify_biblio}
\begin{algorithmic}[1]
\Require Структурированный объект $c = \{\textit{authors}, \textit{title}, \textit{journal}, \textit{year}, \textit{doi}\}$
\Ensure Верифицированный объект $c_{\text{verified}}$ с флагом $V \in \{0, 1\}$

\If{$\textit{doi} \neq \emptyset$}
    \State Отправить запрос к \texttt{api.crossref.org/works/[doi]}
    \If{статус 200 \textbf{и} метаданные совпадают}
        \State \Return $(c, 1)$
    \EndIf
\EndIf

\State Сформировать поисковый запрос $S = \text{``}\textit{authors} + \textit{title} + \textit{year}\text{''}$
\State Выполнить поиск по $S$ в OpenAlex / РИНЦ
\State Вычислить косинусное сходство эмбеддингов и расстояние Левенштейна с топ-3 результатами

\If{$\text{Similarity\_Score} \ge \theta$}
    \State Обогатить $c$ валидным DOI из базы
    \State \Return $(c_{\text{enriched}}, 1)$
\Else
    \State \Return $(c, 0)$ \Comment{Потенциальная галлюцинация}
\EndIf
\end{algorithmic}
\end{algorithm}

Реализация модуля $g_4$ превратит локальный парсер в систему проверки академической добросовестности, защищающую библиометрические базы от загрязнения сгенерированным шумом.